\section{Joint Action--World Modeling}
\label{sec:method}
The design starts from a simple requirement: action prediction and world prediction should be able
to use and update each other's representations. A joint objective specifies which variables are
generated. The attention structure determines how learning those variables interacts.

\subsection{The prediction problem}
Let $o_0$ be the current observation, $\ell$ a language instruction, and $s$ the available
proprioceptive state. We write $h=(o_0,\ell,s)$ for these conditions, with a learned null token for
missing state. Let $z$ denote encoded future observations and $a$ a chunk of future motor actions.
JAM models $p_\theta(z,a\mid h)$. When visual consequence labels $c$ are available, the same
computation models $p_\theta(z,a,c\mid h)$. A fixed video autoencoder maps observations to latents;
its decoder maps generated latents back to images. Action and consequence tokens preserve their
respective control and image-space meanings.

For each generated modality $m\in\{w,a,c\}$, let $x^w=z$, $x^a=a$, and $x^c=c$. We draw independent
standard Gaussian noise $\epsilon^m$ and a corruption level $\tau_m$ between clean and fully noisy.
Rectified-flow training constructs
\begin{equation}
 x^m_{\tau_m}=(1-\tau_m)x^m+\tau_m\epsilon^m,
 \qquad u^m=\epsilon^m-x^m.
 \label{eq:flow}
\end{equation}
Here $u^m$ is the target velocity along the interpolation from data to noise. The model predicts all
available velocities from the corrupted modalities and the clean conditions:
\begin{equation}
 (v^w_\theta,v^a_\theta,v^c_\theta)
 =F_\theta(z_{\tau_w},a_{\tau_a},c_{\tau_c};h,\tau_w,\tau_a,\tau_c).
 \label{eq:prediction}
\end{equation}
The current configuration samples consequence corruption independently of action corruption. This
keeps the confidence of the consequence input distinct from that of the motor input. The formulation
follows independent-modality corruption in UWM \citep{uwm}; the coupling described below specifies
how the pretrained representations participate.

\begin{figure}[t]
\centering
\resizebox{\linewidth}{!}{%
\begin{tikzpicture}[x=1cm,y=1cm,font=\small,
 box/.style={draw=stone,rounded corners=2pt,align=center,line width=.6pt,inner sep=5pt},
 flow/.style={-{Latex[length=2mm]},line width=.8pt,draw=black!65},
 exchange/.style={<->,>=Latex,line width=1.2pt,draw=world!60!black}]
\node[font=\bfseries] at (1.4,6.1) {Demonstrations};
\node[font=\bfseries] at (8.7,6.1) {Joint Action--World Modeling};
\node[box,fill=stone!9,minimum width=8.3cm] (cond) at (8.7,5.45) {Conditions: current observation, instruction, state};
\node[box,fill=stone!10,text width=2.2cm] (robot) at (1.4,4.7) {Robot video\\+ actions};
\node[box,fill=stone!10,text width=2.2cm] (human) at (1.4,3.25) {Human video\\+ tracked motion};
\node[box,fill=stone!6,text width=2.2cm,font=\footnotesize] (prep) at (1.4,1.65) {Encode video\\Map targets\\Apply masks};
\draw[flow] (robot.west) -- (-.05,4.7) -- (-.05,1.65) -- (prep.west);
\draw[flow] (human.south) -- (prep.north);
\node[box,fill=world!10,draw=world,minimum width=2.9cm,minimum height=2.7cm] (w) at (5.3,3.15) {};
\node[box,fill=action!10,draw=action,minimum width=3.4cm,minimum height=2.7cm] (a) at (11.3,3.15) {};
\node[font=\bfseries] at (5.3,4.18) {World stream};
\node[font=\bfseries] at (11.3,4.18) {Agent stream};
\node at (5.3,3.68) {$z_{\tau_w}$};
\node at (11.3,3.68) {state\quad $a_{\tau_a}$\quad $c_{\tau_c}$};
\node[font=\footnotesize,align=center] at (8.3,4.16) {Shared attention\\at every layer};
\foreach \y in {2.35,2.75,3.15} {
 \draw[fill=world!20,draw=world] (4.25,\y) rectangle (6.35,\y+.17);
 \draw[fill=action!22,draw=action] (10.0,\y) rectangle (11.18,\y+.17);
 \draw[fill=consequence!22,draw=consequence] (11.3,\y) rectangle (12.6,\y+.17);
 \draw[exchange] (6.88,\y+.08) -- (9.45,\y+.08);
}
\node[font=\footnotesize,align=center,text width=2.2cm] at (8.2,2.02) {Keys / values\\+ gradients};
\draw[flow] (5.3,5.15) -- (w.north);
\draw[flow] (11.3,5.15) -- (a.north);
\node[font=\footnotesize,fill=white] at (5.3,4.78) {$\tau_w$};
\node[font=\footnotesize,fill=white] at (11.3,4.78) {$\tau_a,\tau_c$};
\node[font=\footnotesize,text=black!60] at (8.3,4.83) {Independent corruption};
\draw[flow] (prep.east) -- (3.3,1.65) -- (3.3,3.15) -- (w.west);
\draw[flow] (3.3,1.65) -- (3.3,1.4) -- (13.5,1.4) -- (13.5,3.15) -- (a.east);
\node[box,fill=world!10,draw=world,minimum width=2.9cm] (vw) at (5.3,.73) {$v^w$\quad Future world};
\node[box,fill=action!10,draw=action,minimum width=2.0cm] (va) at (10.05,.73) {$v^a$\quad Actions};
\node[box,fill=consequence!10,draw=consequence,minimum width=2.35cm] (vc) at (12.7,.73) {$v^c$\quad Consequences};
\draw[flow,draw=world] (w.south) -- (vw.north);
\draw[flow,draw=action] (10.05,1.8) -- (va.north);
\draw[flow,draw=consequence] (12.7,1.8) -- (vc.north);
\node[box,fill=stone!8,minimum width=10.1cm] (loss) at (8.7,-.35) {Joint rectified-flow matching\quad $\mathcal{L}_w+\mathcal{L}_a+\mathcal{L}_c$\quad (valid labels)};
\draw[flow,draw=world] (vw.south) -- (5.3,-.1);
\draw[flow,draw=action] (va.south) -- (10.05,-.1);
\draw[flow,draw=consequence] (vc.south) -- (12.7,-.1);
\end{tikzpicture}}
\caption{\textbf{How action learning enters world pretraining.} Source adapters encode video and map available commands or tracked motion into masked targets. Independently corrupted future latents, actions, and consequences enter coupled world and agent streams. Attention exchanges carry information and gradients throughout the model. Separate velocity readouts receive their corresponding masked targets within the joint loss. World and agent streams share clean observation and language conditions; the state enters through the agent.}
\label{fig:framework}
\end{figure}


\subsection{Reciprocal computation across streams}
The world stream retains a pretrained video transformer's latent embeddings, blocks, and output
projection. The agent stream contains a state token followed by action and consequence tokens. It
uses separate input projections for each token type and noise-dependent adaptive normalization. Both
streams receive the language condition. Their residual widths may differ; their attention
projections share a common head geometry.

At each coupled layer, each stream forms its own queries, keys, and values. World queries attend to
world and agent keys; agent queries attend to agent and world keys. In schematic form, for stream
$m$ and the other stream $\bar m$,
\begin{equation}
 H'_m=H_m+P_m\,\operatorname{Attn}
 \bigl(Q_m,[K_m;K_{\bar m}],[V_m;V_{\bar m}]\bigr),
 \label{eq:coupling}
\end{equation}
where $H_m$ is the layer input, $P_m$ maps attention outputs to its residual width, and brackets
denote token concatenation. Each block then applies language cross-attention and a feed-forward
update. Consequence tokens belong to the agent computation and participate in the same exchange.
Figure~\ref{fig:framework} shows the complete training flow.

The default coupling is bidirectional at every layer, with the exchanged tensors attached to the
computation graph. Thus action loss has a route into world parameters through the world keys and
values read by the agent; world loss has a reciprocal route into agent parameters. These routes
establish a mechanism for action supervision to influence the world representation. Their value for
transfer is an empirical question. Visibility masks and gradient stops expose the mechanism to
controlled comparison.

The current observation forms a clean prefix in the world sequence. Prefix queries attend only to
the clean prefix, while generated world tokens can read the prefix and the agent stream. This
preserves an observation context throughout joint prediction. The state token is a clean input to
the agent stream and is excluded from the generation loss.

\subsection{Learning under partial supervision}
Each example carries coordinate masks $M^m$ for the labels it supplies. We use a masked mean squared
error, $\operatorname{MSE}_{M^m}$, averaged over valid coordinates within that example. Let
$\mathcal{B}_m$ be the batch examples with at least one supervised coordinate for modality $m$. The
implemented objective is
\begin{equation}
 \mathcal{L}=\sum_{m\in\{w,a,c\}}\lambda_m
 \frac{1}{|\mathcal{B}_m|}
 \sum_{i\in\mathcal{B}_m}
 \omega(\tau_{m,i})\,
 \operatorname{MSE}_{M_i^m}\!\left(v^m_{\theta,i},u_i^m\right),
 \label{eq:objective}
\end{equation}
where $\lambda_m$ controls each modality's contribution and $\omega$ weights its noise level. An
empty supervised set contributes zero. World loss covers future latents; the clean prefix is
excluded. Missing coordinates are masked after corruption as well as in the loss. Robot batches
supervise actions and video, while human batches supervise video and the available consequences.
Consequence annotations may also accompany robot data. Appendix~\ref{app:details} gives the
implemented sampling distribution and normalization rules.
